% --- CORRECCION: Forzar interpretacion raw de la entrada ---
\UseRawInputEncoding
\documentclass[aspectratio=169, 10pt, xcolor=table]{beamer}

% --- Configuracion del Tema ---
\usetheme{Madrid}
\usecolortheme{dolphin}

% --- Paquetes necesarios ---
\usepackage[utf8]{inputenc}
\usepackage[T1]{fontenc}
\usepackage[spanish,es-noquoting]{babel}
\usepackage{amsmath, amssymb}
\usepackage{mathtools}
\usepackage{booktabs}
\usepackage{graphicx}
\usepackage{listings}
\usepackage{tikz}
\usetikzlibrary{arrows.meta, positioning, shapes.geometric, babel}
\usepackage{etoolbox}
\usepackage{upquote}

% --- Estilo para codigo SQL y Python ---
\definecolor{codegreen}{rgb}{0,0.6,0}
\definecolor{codegray}{rgb}{0.5,0.5,0.5}
\definecolor{codepurple}{rgb}{0.58,0,0.82}
\definecolor{backcolour}{rgb}{0.95,0.95,0.92}
\lstdefinestyle{mysql}{
    language=SQL,
    backgroundcolor=\color{backcolour},
    commentstyle=\color{codegreen},
    keywordstyle=\color{blue},
    numberstyle=\tiny\color{codegray},
    stringstyle=\color{codepurple},
    basicstyle=\ttfamily\scriptsize,
    breaklines=true,
    frame=single,
    showstringspaces=false,
    literate={á}{{\'a}}1 {é}{{\'e}}1 {í}{{\'i}}1 {ó}{{\'o}}1 {ú}{{\'u}}1 {ñ}{{\~n}}1 {ü}{{\"u}}1
}
\lstdefinestyle{python}{
    language=Python,
    backgroundcolor=\color{backcolour},
    commentstyle=\color{codegreen},
    keywordstyle=\color{blue},
    numberstyle=\tiny\color{codegray},
    stringstyle=\color{codepurple},
    basicstyle=\ttfamily\scriptsize,
    breaklines=true,
    frame=single,
    showstringspaces=false,
    literate={á}{{\'a}}1 {é}{{\'e}}1 {í}{{\'i}}1 {ó}{{\'o}}1 {ú}{{\'u}}1 {ñ}{{\~n}}1 {ü}{{\"u}}1
}
\lstset{style=mysql} % Por defecto SQL

% --- Pie de pagina personalizado ---
\makeatletter
\setbeamertemplate{footline}{
  \leavevmode%
  \hbox{%
  \begin{beamercolorbox}[wd=.333333\paperwidth,ht=2.25ex,dp=1ex,center]{author in head/foot}%
    \usebeamerfont{author in head/foot}\insertshortauthor
  \end{beamercolorbox}%
  \begin{beamercolorbox}[wd=.333333\paperwidth,ht=2.25ex,dp=1ex,center]{title in head/foot}%
    \usebeamerfont{title in head/foot}Sesión 9
  \end{beamercolorbox}%
  \begin{beamercolorbox}[wd=.333333\paperwidth,ht=2.25ex,dp=1ex,right]{date in head/foot}%
    \usebeamerfont{date in head/foot}BBDD \hfill \insertframenumber/\inserttotalframenumber \hspace*{0.3cm}
  \end{beamercolorbox}}%
  \vskip0pt%
}
\makeatother

% --- Datos de la portada ---
\title[SESIÓN 9 BBDD]{\textbf{Sesión 9}}
\subtitle{Data Warehousing, Pipelines y Procesamiento }
\author[Prof. C. Sánchez]{Carlos César Sánchez Coronel}
\institute{\textbf{BASES DE DATOS}}
\date{}

\setbeamertemplate{navigation symbols}{}

\AtBeginSection[]{
  \begin{frame}
    \frametitle{Contenido de la sección}
    \tableofcontents[currentsection]
  \end{frame}
}

\begin{document}

\begin{frame}
    \titlepage
\end{frame}

\begin{frame}{Contenido general}
    \tableofcontents
\end{frame}

% ============================================================================
% SECCIÓN 1: INTRODUCCIÓN
% ============================================================================
\section{Introducción}

\begin{frame}{Objetivos de la sesión}
    \begin{itemize}
        \item Definir y diferenciar Data Warehouse, Data Mart, Data Lake y Lakehouse.
        \item Comprender los procesos ETL y ELT, sus diferencias y cuándo aplicarlos.
        \item Distinguir entre procesamiento batch y streaming (Lambda y Kappa).
        \item Conocer técnicas de carga incremental (watermarks, CDC, MERGE).
        \item Identificar las capas típicas en un data warehouse (raw, staging, core, data marts).
        \item Entender el papel de los metadatos y la gobernanza.
        \item Explorar la integración con herramientas de BI y costos asociados.
        \item Analizar arquitecturas híbridas (SQL + NoSQL).
        \item Definir pipelines, jobs y orquestación (Airflow).
    \end{itemize}
\end{frame}

% ============================================================================
% SECCIÓN 2: FUNDAMENTOS DE ALMACENAMIENTO ANALÍTICO
% ============================================================================
\section{Fundamentos de Almacenamiento Analítico}

\begin{frame}{Data Warehouse (DW)}
    \begin{block}{Características}
        \begin{itemize}
            \item \textbf{Orientado a temas}: ventas, clientes, productos.
            \item \textbf{Integrado}: datos de múltiples fuentes unificados.
            \item \textbf{No volátil}: una vez cargados, no se modifican.
            \item \textbf{Histórico}: series temporales para tendencias.
        \end{itemize}
    \end{block}
    \begin{exampleblock}{Ejemplos}
        Amazon Redshift, Google BigQuery, Snowflake, SQL Server DW.
    \end{exampleblock}
\end{frame}

\begin{frame}{Data Mart}
    \begin{itemize}
        \item Subconjunto del DW orientado a un área de negocio específica (ventas, marketing).
        \item Puede ser \textbf{dependiente} (se alimenta del DW) o \textbf{independiente} (fuente propia).
    \end{itemize}
\end{frame}

\begin{frame}{Data Lake}
    \begin{block}{Repositorio de datos en bruto}
        \begin{itemize}
            \item Almacena cualquier formato: estructurado, semiestructurado, no estructurado.
            \item Bajo costo (S3, HDFS).
            \item \textbf{Schema-on-read}: el esquema se aplica al leer.
            \item Ideal para exploración y ciencia de datos.
        \end{itemize}
    \end{block}
\end{frame}

\begin{frame}{Lakehouse}
    \begin{block}{Fusión de Data Lake y Data Warehouse}
        \begin{itemize}
            \item Datos en formato abierto (Parquet).
            \item Transacciones ACID, evolución de esquema.
            \item Rendimiento analítico similar al DW.
        \end{itemize}
    \end{block}
    \begin{exampleblock}{Tecnologías}
        Delta Lake (Databricks), Apache Iceberg, Apache Hudi.
    \end{exampleblock}
    \begin{table}
        \centering
        \begin{tabular}{l|c|c|c}
            \toprule
            Característica & DW & Data Lake & Lakehouse \\
            \midrule
            Formato de datos & Estructurado & Cualquier formato & Abierto (Parquet) \\
            Esquema & Fijo (schema-on-write) & Flexible (schema-on-read) & Evolutivo \\
            Transacciones ACID & Sí & Limitado & Sí \\
            Rendimiento analítico & Alto & Bajo/Medio & Alto \\
            Costo & Alto & Bajo & Medio \\
            \bottomrule
        \end{tabular}
    \end{table}
\end{frame}

% ============================================================================
% SECCIÓN 3: ETL vs ELT
% ============================================================================
\section{Procesos ETL vs ELT}

\begin{frame}{ETL (Extract, Transform, Load)}
    \begin{center}
        \begin{tikzpicture}[node distance=2cm, auto]
            \node (source) [draw, rectangle] {Fuentes};
            \node (staging) [draw, rectangle, right of=source, xshift=1.5cm] {Staging (Transformación)};
            \node (dw) [draw, rectangle, right of=staging, xshift=1.5cm] {Data Warehouse};
            \draw[->] (source) -- node[above] {Extract} (staging);
            \draw[->] (staging) -- node[above] {Load} (dw);
        \end{tikzpicture}
    \end{center}
    \begin{itemize}
        \item \textbf{Ventajas}: control total, menor carga en el DW, adecuado para datos sensibles.
        \item \textbf{Desventajas}: infraestructura adicional, mayor latencia.
    \end{itemize}
\end{frame}

\begin{frame}{ELT (Extract, Load, Transform)}
    \begin{center}
        \begin{tikzpicture}[node distance=2cm, auto]
            \node (source) [draw, rectangle] {Fuentes};
            \node (dw) [draw, rectangle, right of=source, xshift=2cm] {Data Warehouse};
            \draw[->] (source) -- node[above] {Extract \& Load} (dw);
            \node (transform) [draw, rectangle, below of=dw, yshift=-1cm] {Transformaciones (SQL)};
            \draw[->] (dw) -- (transform);
        \end{tikzpicture}
    \end{center}
    \begin{itemize}
        \item \textbf{Ventajas}: simplicidad, aprovecha escalabilidad del DW, datos crudos disponibles rápidamente.
        \item \textbf{Desventajas}: costo de cómputo en el DW, no ideal para no estructurados.
    \end{itemize}
\end{frame}

\begin{frame}{¿Cuándo usar cada uno?}
    \begin{block}{ETL}
        \begin{itemize}
            \item Transformaciones muy complejas o datos sensibles.
            \item Entornos on-premise.
            \item Herramientas: Informatica, Talend, Spark batch.
        \end{itemize}
    \end{block}
    \begin{block}{ELT}
        \begin{itemize}
            \item DW en la nube (Snowflake, BigQuery, Redshift).
            \item Agilidad y grandes volúmenes.
            \item Herramientas: dbt, SQL directo.
        \end{itemize}
    \end{block}
\end{frame}

% ============================================================================
% SECCIÓN 4: PROCESAMIENTO BATCH VS STREAMING
% ============================================================================
\section{Procesamiento Batch vs Streaming}

\begin{frame}{Batch (Lotes)}
    \begin{itemize}
        \item Datos acumulados durante un período.
        \item Ejecución programada (diaria, horaria).
        \item Herramientas: Apache Spark, Hadoop MapReduce, scripts con cron.
        \item Casos: facturación mensual, informes de cierre.
    \end{itemize}
\end{frame}

\begin{frame}{Streaming (Tiempo Real)}
    \begin{itemize}
        \item Procesamiento continuo de eventos.
        \item Latencias de milisegundos a segundos.
        \item Herramientas: Apache Kafka, Apache Flink, Spark Structured Streaming.
        \item Casos: detección de fraude, monitoreo IoT.
    \end{itemize}
\end{frame}

\begin{frame}{Arquitecturas Lambda y Kappa}
    \begin{columns}[t]
        \column{0.5\textwidth}
        \textbf{Lambda}
        \begin{itemize}
            \item Dos capas: batch y speed.
            \item Precisión batch + baja latencia.
            \item Compleja de mantener.
        \end{itemize}
        \column{0.5\textwidth}
        \textbf{Kappa}
        \begin{itemize}
            \item Una sola capa de streaming.
            \item Reprocesamiento como streaming.
            \item Más simple (requiere motor potente).
        \end{itemize}
    \end{columns}
    \vspace{0.5cm}
    \begin{center}
        \begin{tikzpicture}[node distance=1.5cm, auto]
            \node (source) [draw] {Fuente};
            \node (batch) [draw, below left of=source, xshift=-1cm] {Capa batch};
            \node (speed) [draw, below right of=source, xshift=1cm] {Capa speed};
            \node (service) [draw, below of=source, yshift=-2cm] {Capa de servicio};
            \draw[->] (source) -- (batch);
            \draw[->] (source) -- (speed);
            \draw[->] (batch) -- (service);
            \draw[->] (speed) -- (service);
            \node at (source.north) [above] {Lambda};
        \end{tikzpicture}
    \end{center}
\end{frame}

% ============================================================================
% SECCIÓN 5: TÉCNICAS DE CARGA INCREMENTAL
% ============================================================================
\section{Técnicas de Carga Incremental}

\begin{frame}{Marcas de agua (Watermarks)}
    \begin{itemize}
        \item Guardar la fecha/hora de la última carga.
        \item Extraer registros posteriores a esa marca.
    \end{itemize}
    \begin{block}{Ejemplo SQL}
        \begin{lstlisting}
-- Obtener última fecha cargada
SELECT MAX(fecha) INTO ultima_fecha FROM dw_ventas;

-- Insertar nuevos datos
INSERT INTO dw_ventas
SELECT * FROM stg_ventas WHERE fecha > ultima_fecha;
        \end{lstlisting}
    \end{block}
\end{frame}

\begin{frame}{Change Data Capture (CDC)}
    \begin{itemize}
        \item Captura cambios en tiempo real desde logs de BD.
        \item Detecta inserciones, actualizaciones, eliminaciones.
        \item Herramientas: Debezium, AWS DMS, Oracle GoldenGate.
    \end{itemize}
\end{frame}

\begin{frame}[fragile]{MERGE (UPSERT)}
    \begin{block}{PostgreSQL / SQL Server}
        \begin{lstlisting}
MERGE INTO dw_clientes AS target
USING stg_clientes AS source
ON target.id = source.id
WHEN MATCHED THEN
    UPDATE SET nombre = source.nombre, updated_at = CURRENT_TIMESTAMP
WHEN NOT MATCHED THEN
    INSERT (id, nombre, created_at)
    VALUES (source.id, source.nombre, CURRENT_TIMESTAMP);
        \end{lstlisting}
    \end{block}
\end{frame}

% ============================================================================
% SECCIÓN 6: CAPAS EN UN DATA WAREHOUSE
% ============================================================================
\section{Capas en un Data Warehouse}

\begin{frame}{Capas típicas}
    \begin{description}
        \item[Raw (Bronze)] Datos crudos tal como llegan (CSV, JSON). Sirve para reprocesar.
        \item[Staging (Silver)] Datos limpios, validados, con tipos correctos, sin duplicados.
        \item[Core / Integración (Gold)] Modelo dimensional (hechos y dimensiones) listo para negocio.
        \item[Data Marts] Subconjuntos para áreas específicas (ventas, marketing).
    \end{description}
\end{frame}

\begin{frame}{Ejemplo de diseño}
    \begin{itemize}
        \item \textbf{Raw}: archivos JSON de ventas diarias, CSV de productos, logs web.
        \item \textbf{Staging}: tablas temporales con ventas limpias, clientes unificados.
        \item \textbf{Core}: tabla de hechos \texttt{ventas}, dimensiones \texttt{cliente}, \texttt{producto}, \texttt{tiempo}.
        \item \textbf{Data Marts}: \texttt{ventas\_por\_region}, \texttt{rentabilidad\_producto}.
    \end{itemize}
\end{frame}

% ============================================================================
% SECCIÓN 7: METADATOS Y GOBERNANZA
% ============================================================================
\section{Metadatos y Gobernanza}

\begin{frame}{Tipos de metadatos}
    \begin{description}
        \item[Técnicos] Esquemas, tipos de datos, particiones, estadísticas.
        \item[De negocio] Definiciones, reglas de cálculo, propietarios.
        \item[Operacionales] Logs de ejecución, tiempos de carga.
    \end{description}
\end{frame}

\begin{frame}{Herramientas de catálogo y linaje}
    \begin{itemize}
        \item Catálogos: Apache Atlas, Amundsen, DataHub, AWS Glue Catalog.
        \item Linaje: trazabilidad del origen y transformaciones (importante para auditoría y GDPR).
    \end{itemize}
\end{frame}

% ============================================================================
% SECCIÓN 8: INTEGRACIÓN CON HERRAMIENTAS DE BI
% ============================================================================
\section{Integración con BI}

\begin{frame}{Conexión a herramientas de visualización}
    \begin{itemize}
        \item Power BI, Tableau, Looker, Qlik.
        \item Modos: \textbf{Importación} (datos copiados) vs \textbf{DirectQuery} (consultas en vivo).
        \item Recomendación: usar agregaciones precalculadas y vistas materializadas para rendimiento.
    \end{itemize}
\end{frame}

% ============================================================================
% SECCIÓN 9: DIMENSIONAMIENTO Y COSTOS
% ============================================================================
\section{Dimensionamiento de Almacenamiento}

\begin{frame}{Factores a considerar}
    \begin{itemize}
        \item Volumen diario de datos.
        \item Retención histórica.
        \item Compresión (columnar: Parquet, ORC).
        \item Índices y particionamiento.
        \item Backups y replicación.
    \end{itemize}
\end{frame}

\begin{frame}{Ejemplo de cálculo}
    \textbf{Problema:} 500 GB/día, retención 3 años, compresión 4:1, costo S3 \$0.023/GB-mes.

    \begin{itemize}
        \item Datos diarios: 500 GB.
        \item Anual: 500 GB * 365 = 182.5 TB.
        \item 3 años: 547.5 TB.
        \item Con compresión 4:1: 136.9 TB.
        \item Margen 10\%: 150.6 TB.
        \item Costo mensual: 150.6 * 1024 * 0.023 ≈ \$3,546/mes.
    \end{itemize}
\end{frame}

% ============================================================================
% SECCIÓN 10: ARQUITECTURAS HÍBRIDAS (SQL + NoSQL)
% ============================================================================
\section{Arquitecturas Híbridas}

\begin{frame}{Caso de uso combinado}
    \begin{itemize}
        \item \textbf{PostgreSQL}: transacciones (pedidos).
        \item \textbf{MongoDB}: catálogo de productos flexible.
        \item \textbf{Elasticsearch}: búsqueda rápida.
        \item \textbf{Snowflake}: análisis de ventas.
    \end{itemize}
\end{frame}

\begin{frame}{Sincronización entre motores}
    \begin{itemize}
        \item CDC con Debezium desde PostgreSQL a Kafka.
        \item Consumidores escriben en MongoDB, Elasticsearch y el DW.
        \item Batch periódico con Apache Spark para reconciliación.
    \end{itemize}
\end{frame}

% ============================================================================
% SECCIÓN 11: PIPELINES, JOBS Y ORQUESTACIÓN
% ============================================================================
\section{Pipelines, Jobs y Orquestación}

\begin{frame}{Definiciones}
    \begin{description}
        \item[Pipeline] Flujo completo de procesamiento de datos.
        \item[Job] Unidad de trabajo (puede contener varias tareas).
        \item[Tarea] Paso individual (script, consulta SQL).
        \item[DAG] Grafo acíclico dirigido que define dependencias.
    \end{description}
\end{frame}

\begin{frame}{Orquestadores populares}
    \begin{itemize}
        \item Apache Airflow (el más usado, DAGs en Python).
        \item Prefect, Dagster (alternativas modernas).
        \item AWS Step Functions, Google Cloud Composer (gestionados).
    \end{itemize}
\end{frame}

\begin{frame}[fragile]{Ejemplo de DAG en Airflow}
    \begin{lstlisting}[style=python]
from airflow import DAG
from airflow.operators.python import PythonOperator
from airflow.providers.apache.spark.operators.spark_submit import SparkSubmitOperator
from airflow.providers.amazon.aws.operators.redshift import RedshiftSQLOperator
from datetime import datetime
import requests

def extraer_api():
    response = requests.get('https://api.ejemplo.com/data')
    with open('/tmp/data.json', 'w') as f:
        f.write(response.text)

with DAG('pipeline_ventas', start_date=datetime(2025,1,1), schedule='@daily') as dag:
    extraer = PythonOperator(task_id='extraer', python_callable=extraer_api)
    transformar = SparkSubmitOperator(
        task_id='transformar',
        application='/path/to/spark_job.py',
        conn_id='spark_default'
    )
    cargar = RedshiftSQLOperator(
        task_id='cargar',
        sql='CALL cargar_ventas();',
        redshift_conn_id='redshift_default'
    )
    extraer >> transformar >> cargar
    \end{lstlisting}
\end{frame}

% ============================================================================
% SECCIÓN 12: EJERCICIOS RESUELTOS
% ============================================================================
\section{Ejercicios Resueltos}

\begin{frame}{Ejercicio 1: Diseño de capas para un DW de ventas}
    \textbf{Enunciado:} Proponga un diseño de capas (raw, staging, core) y las transformaciones típicas.
    \begin{block}{Solución}
        \begin{itemize}
            \item \textbf{Raw}: Archivos CSV/JSON de ventas, productos, clientes tal como llegan. Particionado por fecha.
            \item \textbf{Staging}: Asignar tipos correctos, eliminar duplicados, unificar formatos de fecha. Guardar en Parquet.
            \item \textbf{Core}: Construir tablas de hechos (ventas) y dimensiones (cliente, producto, tiempo) con modelado dimensional. Aplicar SCD tipo 2.
        \end{itemize}
    \end{block}
\end{frame}

\begin{frame}[fragile]{Ejercicio 2: Carga incremental con MERGE}
    \textbf{Enunciado:} Escribir una consulta SQL que actualice \texttt{dw\_ventas} con los datos de \texttt{stg\_ventas} del día usando MERGE.
    \begin{block}{Solución}
        \begin{lstlisting}
MERGE INTO dw_ventas AS target
USING stg_ventas AS source
ON target.id = source.id
WHEN MATCHED THEN
    UPDATE SET target.monto = source.monto,
               target.fecha_actualizacion = CURRENT_TIMESTAMP
WHEN NOT MATCHED THEN
    INSERT (id, fecha, monto, fecha_actualizacion)
    VALUES (source.id, source.fecha, source.monto, CURRENT_TIMESTAMP);
        \end{lstlisting}
    \end{block}
\end{frame}

\begin{frame}{Ejercicio 3: Estimar costo de almacenamiento}
    \textbf{Enunciado:} 500 GB/día, retención 3 años, compresión 4:1. Calcular almacenamiento y costo mensual en S3 (\$0.023/GB-mes).
    \begin{block}{Solución}
        \begin{itemize}
            \item Datos diarios: 500 GB.
            \item Anual: 500 * 365 = 182.5 TB.
            \item 3 años: 547.5 TB.
            \item Compresión 4:1: 136.9 TB.
            \item Margen 10\%: 150.6 TB.
            \item Costo mensual: 150.6 * 1024 * 0.023 ≈ \$3,546/mes.
        \end{itemize}
    \end{block}
\end{frame}

\begin{frame}[fragile]{Ejercicio 4: Pipeline con Airflow}
    \textbf{Enunciado:} Crear DAG con tareas: extraer de API, transformar con Spark, cargar en Redshift.
    \begin{block}{Solución (ver código en sección 11)}
        (El código del DAG se mostró anteriormente en la diapositiva de Airflow).
    \end{block}
\end{frame}

% ============================================================================
% SECCIÓN 13: GLOSARIO
% ============================================================================
\section{Glosario}

\begin{frame}{Términos Clave}
    \begin{description}
        \item[Batch] Procesamiento por lotes en intervalos.
        \item[CDC] Change Data Capture, captura de cambios en tiempo real.
        \item[Data Lake] Repositorio de datos en bruto.
        \item[Data Mart] Subconjunto del DW para un área específica.
        \item[Data Warehouse] Almacén de datos optimizado para análisis.
        \item[DAG] Directed Acyclic Graph, dependencias de tareas.
        \item[ELT] Extract, Load, Transform (cargar primero, transformar después).
        \item[ETL] Extract, Transform, Load (transformar antes de cargar).
        \item[Lakehouse] Fusión de data lake y data warehouse.
        \item[Linaje] Trazabilidad del origen y transformaciones.
        \item[Merge] Operación que inserta o actualiza según exista.
        \item[Metadatos] Datos sobre los datos.
        \item[Orquestación] Coordinación de tareas en un pipeline.
        \item[Pipeline] Flujo completo de procesamiento.
        \item[Streaming] Procesamiento continuo en tiempo real.
    \end{description}
\end{frame}

% --- Diapositiva final ---
\begin{frame}
    \centering

    
    \vspace{1cm}
    \large ¡Gracias!
\end{frame}

\end{document}